% --- Archon physics formalization source begin ---
% archon:physics
% archon:covers IPhO2026Problems/problem_ipho_2026_t2_b3.lean
% archon:source-report reports/ipho_2026/problem_ipho_2026_t2_b3.source.json
% archon:problem-id ipho_2026_t2
% archon:part-id T2-B3
% archon:previous-part ipho_2026_t2_b1
% archon:previous-part-policy derive_inline_from_problem_only_material
% archon:previous-part ipho_2026_t2_b2
% archon:previous-part-policy derive_inline_from_problem_only_material

\chapter{Ipho Physics problem ipho\_2026\_t2\_b3}
\label{ch:IPhO2026Problems_problem_ipho_2026_t2_b3}

\paragraph{Problem source.}
\#\# Physical scenario

A half-cylindrical mirror of radius R illuminates a fully absorbing cylindrical
container of radius a.  Their axes are parallel, and the container center lies
R/2 from the mirror center on the symmetry plane.  Uniform parallel sunlight
arrives along the optical axis.  Any ray absorbed by the container reflects at
most once.  Let theta\_max be the largest incidence angle on the mirror among
rays that strike the container, and let P\_0 be the power the cylinder would
receive without the mirror.  See Figure 2f.

\#\# Current subquestion T2-B3

For R = 1.0 m, find a such that P = 5*P\_0, and report it in cm.

\paragraph{Shared problem context.}
A half-cylindrical mirror of radius R illuminates a fully absorbing cylindrical
container of radius a.  Their axes are parallel, and the container center lies
R/2 from the mirror center on the symmetry plane.  Uniform parallel sunlight
arrives along the optical axis.  Any ray absorbed by the container reflects at
most once.  Let theta\_max be the largest incidence angle on the mirror among
rays that strike the container, and let P\_0 be the power the cylinder would
receive without the mirror.  See Figure 2f.

\paragraph{Current subquestion.}
For R = 1.0 m, find a such that P = 5*P\_0, and report it in cm.

\paragraph{Figure/image path.}
ipho\_2026\_source/image/T2\_page-3.png

\paragraph{Reusable previous-part conclusions.}
\begin{itemize}
\item Source T2-B1. Question: Given a = alpha*sin(theta\_max) + beta*sin(2*theta\_max), determine alpha and beta in terms of R. Policy: derive\_inline\_from\_problem\_only\_material
\item Source T2-B2. Question: Express the power ratio P/P\_0 in terms of theta\_max. Policy: derive\_inline\_from\_problem\_only\_material
\end{itemize}

\paragraph{Formalization target.}
create a compiling Lean file with sorry bodies at `IPhO2026Problems/problem\_ipho\_2026\_t2\_b3.lean`.
The Lean declarations must preserve the physical quantities, dimensions or dimensional roles, figure labels, governing-law hypotheses, and final relation expressed by this problem.
Use Mathlib/Physlib names found through LeanExplore where available. If a domain API is missing, introduce faithful local abstractions rather than scalar placeholder aliases.
This is an answer-blind solve-phase task. Derive a candidate result only from the problem-side evidence above; do not assume a target value, select a tolerance around a desired value, or encode a desired result into a candidate domain.
For numerical questions, define the raw end-to-end quantity and apply an explicit source-derived rounding rule. For identification questions, characterize or prove uniqueness before recording the derived candidate.

\begin{theorem}[Ipho Physics formalization target]
\label{thm:physics:ipho_2026_t2_b3:target}
\lean{IPhO2026.T2B3.container_radius_of_power_five, IPhO2026.T2B3.t2_b3_container_radius_cm}
\leanok
The assigned autoformalize agent should translate this IPhO physics problem into Lean declarations in the covered file, with theorem and lemma proof bodies written as `by sorry`.
\end{theorem}
\begin{proof}
\leanok
This is an autoformalization task, not a proof task. Produce faithful statements that can later be proved without weakening the source contract.

\emph{Prover-stage completion.} All 40 proof obligations in
\texttt{IPhO2026Problems/problem\_ipho\_2026\_t2\_b3.lean} are now closed with
kernel-checked proofs (no \texttt{sorry}; only the standard axioms
\texttt{propext}, \texttt{Classical.choice}, \texttt{Quot.sound}).  The proof
chain: the shared T2-B ray-optics model (specular reflection
\texttt{specularReflect} in the radial normal, reflected direction
$d' = (-\sin 2\theta, -\cos 2\theta)$ in \texttt{reflectedDir\_eq},
distance identity \texttt{distToLine\_container}, strike characterization
\texttt{strikes\_iff\_dist\_le}, grazing contact \texttt{dist\_at\_max\_eq} via the
intermediate value theorem and strict monotonicity \texttt{dist\_strictMonoOn})
gives the T2-B1 relation \texttt{container\_radius\_eq\_factored}
($a = R \sin\theta_{\max}(1 - \cos\theta_{\max})$) and, through the absorbed
impact set \texttt{absorbedImpactSet\_eq\_Icc} $= [-R\sin\theta_{\max},
R\sin\theta_{\max}]$, the T2-B2 ratio \texttt{power\_ratio}
($P/P_0 = 1/(1-\cos\theta_{\max})$).  Imposing $P = 5 P_0$ forces
$\cos\theta_{\max} = 4/5$ (\texttt{cos\_thetaMax\_of\_power\_five}),
$\sin\theta_{\max} = 3/5$ on the geometric branch $0 < \theta_{\max} \le \pi/2$
(\texttt{sin\_thetaMax\_of\_power\_five}; uniqueness
\texttt{thetaMax\_unique\_of\_cos\_eq}), hence
$a = 3R/25$ (\texttt{container\_radius\_of\_power\_five}), and at $R = 1$ m the
exact unit conversion \texttt{mToCm} yields $a = 12$ cm
(\texttt{t2\_b3\_container\_radius\_cm}).
\end{proof}
% --- Archon physics formalization source end ---
